\documentclass[a4paper,10pt]{book}
\usepackage[utf8]{inputenc}
\usepackage[T1]{fontenc}
\usepackage[italian]{babel}
\usepackage{pdfpages}
\usepackage{tikz}
\usepackage{eso-pic}
\usepackage[top=3cm, bottom=3cm, left=3cm, right=3cm]{geometry}

\usepackage{xcolor}
\usepackage[most]{tcolorbox}
\usepackage{fontawesome5} % Per inserire le icone (es. \faGithub)


% Definizione della palette sequenziale (Viridis-like)
\definecolor{colorA}{HTML}{003f5c} % 1 - Blu Profondo
\definecolor{colorB}{HTML}{2f4b7c} % 2 - Indaco
\definecolor{colorC}{HTML}{665191} % 3 - Viola Bluastro
\definecolor{colorD}{HTML}{a05195} % 4 - Magenta Scuro
\definecolor{colorE}{HTML}{d45087} % 5 - Rosso Porpora
\definecolor{colorF}{HTML}{f95d6a} % 6 - Arancio-Rosso
\definecolor{colorG}{HTML}{ff7c43} % 7 - Arancio Brillante

\usepackage{hyperref} % Per link interni (es. tabella dei contenuti, riferimenti)
\hypersetup{
    colorlinks=true,
    linkcolor=colorA,
    filecolor=colorE,      
    urlcolor=colorB,
}


% --- Copertina Bio-Fisica (Sfondo con Gradiente Elegante) ---
\begin{titlepage}
    \thispagestyle{empty}
    
    \begin{tikzpicture}[remember picture, overlay]
        
        % 1. Sfondo con gradiente: dal Blu Profondo (colorA) all'Indaco (colorB)
        \shade[top color=colorA, bottom color=colorC] (current page.north west) rectangle (current page.south east);
        
        % 2. Fascia laterale a contrasto
        \fill[colorE] (current page.south west) rectangle ([xshift=0.4cm]current page.north west);
        
        % 3. Rete di Hopfield (Pura, solo linee bianche e nodi)
        \begin{scope}[shift={([xshift=-1.5cm, yshift=-8cm]current page.east)}]
            \def\numnodes{24} 
            \def\radius{9} 
            
            % Disegna tutte le connessioni in bianco semitrasparente
            \foreach \i in {1,...,\numnodes} {
                \foreach \j in {\i,...,\numnodes} {
                    \draw[white, opacity=0.15, line width=0.3pt] ({\i*360/\numnodes}:\radius) -- ({\j*360/\numnodes}:\radius);
                }
            }
            
            % Disegna i nodi con un leggero alone
            \foreach \i in {1,...,\numnodes} {
                % Alone luminoso morbido
                \node[circle, fill=white, opacity=0.15, inner sep=0pt, minimum size=14pt] at ({\i*360/\numnodes}:\radius) {};
                % Core del nodo
                \node[circle, fill=white, inner sep=0pt, minimum size=4pt] at ({\i*360/\numnodes}:\radius) {};
            }
        \end{scope}
    \end{tikzpicture}

    % --- Area Testi: Intestazione ---
    \begin{flushleft}
        \vspace*{1.5cm} % Margine dall'alto
        \textcolor{colorE}{\rule{2cm}{2pt}}\\[0.4cm] % Linea decorativa
        \textcolor{white}{\sffamily\Large\textbf{SAPIENZA UNIVERSITÀ DI ROMA}}\\[0.2cm]
        \textcolor{white!70}{\sffamily\large Laurea Magistrale in Fisica}\\
        
        \vspace{4.5cm}
        
        % Titolo Principale 
        \textcolor{colorE}{\sffamily\fontsize{60}{70}\selectfont\textbf{Laboratorio di}}\\[0.1cm]
        \textcolor{colorE}{\sffamily\fontsize{60}{70}\selectfont\textbf{Biosistemi I}}\\[0.6cm]
        \textcolor{white!80}{\sffamily\Large Appunti del corso}\\[0.5cm]
        \textcolor{white!30}{\rule{6cm}{1pt}} % Sottile linea di separazione
    \end{flushleft}
    
    \vfill % Spinge le info dell'autore verso il fondo
    
    % --- Area Testi: Autore ---
    \begin{flushleft}
        \textcolor{white!60}{\sffamily\large A cura di}\\[0.2cm]
        \textcolor{white}{\sffamily\huge\textbf{Elisabetta Agnello}}\\[1.2cm]
        
        \textcolor{white!60}{\sffamily\large Anno Accademico}\\[0.2cm]
        \textcolor{white}{\sffamily\Large\textbf{2024/2025}}
    \end{flushleft}
    
    \vspace*{1cm} % Margine dal basso
\end{titlepage}
% --- Fine Copertina ---
    

% --- WATERMARK ---
\newcommand{\WatermarkOverlay}{%
  \begin{tikzpicture}[remember picture, overlay]
    \node[rotate=45, scale=3, opacity=0.15, text=black!50] at (current page.center) {
      \sffamily\bfseries github.com / elisabettagnello
    };
  \end{tikzpicture}%
}

\begin{document}

\maketitle

% --- Disclaimer ---
\chapter*{Disclaimer}
\thispagestyle{empty} % Rimuove header e footer solo per questa pagina

Il presente documento costituisce una raccolta di appunti personali presi durante le lezioni del corso di \textit{Laboratorio di Biosistemi I}. 

Si prega di notare che questo materiale \textbf{non è stato sottoposto a una revisione sistematica} né è stato verificato dal docente. Pertanto, è possibile che siano presenti:
\begin{itemize}
    \item Errori concettuali o imprecisioni tecniche;
    \item Refusi o errori di battitura;
    \item Parti incomplete o sezioni mancanti.
\end{itemize}

Questi appunti sono destinati esclusivamente come supporto allo studio e non sostituiscono in alcun modo il materiale didattico ufficiale. L'autore non si assume alcuna responsabilità per l'uso delle informazioni qui contenute.


% =============================================================================
\vspace{3cm} % Spazio verticale per staccare il box dal testo sopra

\begin{tcolorbox}[
    colback=colorA!5,      
    colframe=colorA,       
    title={Contatti e Segnalazioni},
    fonttitle=\bfseries\sffamily,
    arc=3mm,               
    boxrule=1.2pt,         
    halign title=flush left
]
Per consultare l'archivio aggiornato, segnalare refusi o proporre integrazioni agli appunti, puoi utilizzare i seguenti canali:
\begin{description}
    \item[\textsf{\faGlobe\ Archivio Appunti:}]
    \href{https://elisabettagnello.github.io/elisabetta-physics-notes/}{\texttt{elisabettagnello.github.io/elisabetta-physics-notes/}}
    \item[\textsf{\faGithub\ GitHub:}] \href{https://github.com/elisabettagnello}{\texttt{github.com/elisabettagnello}}
    \item[\textsf{\faLinkedin\ LinkedIn:}] \href{https://www.linkedin.com/in/elisabetta-agnello}{\texttt{linkedin.com/in/elisabetta-agnello}}
    \item[\textsf{\faEnvelope\ Email:}] \href{mailto:agnelloe24@gmail.com}{\texttt{agnelloe24@gmail.com}}
\end{description}
\end{tcolorbox}
% =============================================================================


\clearpage % Forza il passaggio alla pagina successiva

\includepdf[pages=-, fitpaper=true, pagecommand={\thispagestyle{empty} \WatermarkOverlay}]{01.pdf}
\includepdf[pages=-, fitpaper=true, pagecommand={\thispagestyle{empty} \WatermarkOverlay}]{02.pdf}


\end{document}
